\section{Introduction}
\label{sec:intro}

The grid runs on old protocols, and an attacker learns them before it breaks them. Power utilities
move measurements and control commands over protocols such as DNP3 (Distributed Network
Protocol~3)~\cite{dnp3std}, which, like most of its peers, ships with no encryption and no
authentication by default. Anyone on the path can read the exchange. The cost of that exposure is not
hypothetical: in December~2015 remote intruders reached into a Ukrainian power grid and left
225{,}000 residents in the dark~\cite{ukraine2015}. An attack at that scale does not begin with the
attack. It begins with watching, with learning which devices are on the wire and how each behaves,
because a control-altering step such as a false-data-injection attack on state estimation only works
once the attacker knows the target~\cite{liufdia2009,cardenas2011}. What a quiet observer can pull out
of unencrypted control traffic is therefore worth taking seriously.

Reading the payload is not even necessary. Traffic-analysis work has shown that packet sizes and
timing on their own recover the content of even encrypted sessions~\cite{sirinam2018,feghhi2016}, and
that a device or a device type can be told apart from its wire-side behavior
alone~\cite{kohno2005,gtid2015}. A useful distinction runs through this line of work: an observer may
seek a device's \emph{identity}, its \emph{type}, or the \emph{class of transaction} it is currently
running. Our concern is the last two: features that let an observer separate one device's transaction
classes and, across devices, its type, without decoding anything.

Two features carry that signature for DNP3, and both are visible on unencrypted traffic even though
the payload is fully parseable. The first is timing. Formby et~al.\ named the cross-layer response
time (CLRT) as a device fingerprint for industrial control systems~\cite{formby2016}: the gap between
the transport-layer acknowledgment and the application-layer response reflects the device's internal
processing and stays stable across sessions. The second is shape. A DNP3 response is delivered in TCP
segments, and the vector of segment sizes an observer records is itself a fingerprint, because the
outstation lays out its frames deterministically and so produces the same shape every time.

General defenses against such leakage already exist, and they reshape observable traffic in
well-understood ways. Padding hides a flow's size by adding bytes; morphing transforms one class's
size distribution into another's~\cite{wright2009}; adaptive padding breaks the inter-packet timing
pattern~\cite{juarez2016}; constant-rate and burst-molding schemes bound the leak at a fixed
cost~\cite{tamaraw2014,walkietalkie2017}; differential privacy quantifies the residual
leak~\cite{netshaper2024}; and recent systems push shaping into the network on programmable
switches~\cite{ditto2022,securitas2026,minos2025}. Each of these carries an assumption about where it
runs and what it may change.

Those assumptions do not transfer cleanly to DNP3. A control center periodically polls a field
device, the two endpoints run fixed vendor firmware that an operator cannot modify, the devices are
long-lived and repeatedly polled, and the protocol is validated on receipt: DNP3 protects every
16-byte block of a message with a cyclic redundancy check (CRC) that the master verifies as it
reassembles. A defense that pads or rewrites the response bytes therefore risks a CRC failure the
master rejects, a defense that changes the endpoints is unavailable, and a defense that reshapes only
size or only timing leaves the other feature intact.

Consider one concrete read on the relay we evaluate. The master sends a READ request; the outstation
returns first a TCP segment carrying the acknowledgment flag and then, a moment later, its DNP3
application response. The CLRT this paper normalizes is the interval between exactly those two
packets: the first acknowledgment-bearing segment and the first DNP3 application response
(function~0x81). On our SEL-751 relay that gap has a median near 1.3~ms for a read and near 2.1~ms for
a control select, and the response is a single 49-byte segment, a shape we write~$[49]$. Those two
numbers alone separate the read and select transaction classes, which is the transaction-class leak
Formby's fingerprint predicts for this device.

The requirements for closing this leak follow from the constraints rather than from taste. The
defense may not add or rewrite DNP3 bytes, so it cannot pad; it may not change the endpoints, so it
cannot run on the master or the outstation; a controller-mediated release would add its own timing
variability and a deployment dependency; size-only shaping would leave the timing feature and
timing-only shaping would leave the shape feature; and whatever it does must still deliver a valid,
ordered DNP3 exchange the master accepts. The one place that satisfies all of these is a programmable
switch on the path, which can retime and re-segment the traffic without touching either endpoint.

We place the defense in the data plane of one programmable switch at the outstation's edge. It does
two things. For a response, it holds the acknowledgment and the response and releases each at a fixed
offset from the request, so the master-visible CLRT becomes one policy value; and it splits the
49-byte response into a fixed two-segment shape, $[28,21]$, cutting only at a boundary the protocol
already places between CRC-protected blocks, so no DNP3 byte changes and the master reassembles a
valid 49-byte response. A control command needs a second treatment, because an outbound OPERATE
creates a different causal problem: the delay from the OPERATE to its confirmation is itself an
operation-timing fingerprint, and simply releasing the OPERATE on a fixed schedule would still expose
it. The defense instead holds the OPERATE and releases it once at an internal delay drawn from a
bounded set, while pinning the master-visible acknowledgment and echo to the request's arrival, so the
gap the observer measures does not vary with that delay.

Building this on real hardware raised systems problems that shaped the design. The switch has no
general-purpose timer, so a delayed release must be realized from packet scheduling. Several events
must be anchored to a single request timestamp rather than to whatever arrives next. A valid response
must be split without any oracle that would let us claim byte identity to a source frame. The two
treatments cannot share one priority schedule without the response's queues starving the command's,
so they must run on separate internal scheduling lanes. The whole program must fit one switch
pipeline. And the mechanism must operate safely on a physical protective relay.

We evaluate the design on a testbed with a physical master, one Tofino-1 switch, and a physical
SEL-751 relay, comparing native and defended traffic captured on the same loaded program under a
runtime mode change. On the master-facing link the defended CLRT is a fixed 4.001~ms for both read and
select in place of the native 1.3 and 2.1~ms; every one of 1{,}280 defended eligible responses
presents the $[28,21]$ shape with a sequence-contiguous, CRC-valid, checksum-valid 49-byte
reconstruction and no unsplit escape; the control-command gap stays at 4.00~ms across the evaluated
internal delays; and a read-versus-select timing classifier drops from 0.59 to the 0.50 chance
baseline. The evaluation is bounded: we test one outstation and one eligible response class, so we
replace that device's evaluated signature rather than establishing indistinguishability among
devices, and we do not observe the relay-facing link, so the single internal release is verified by
the implementation model and not measured on the wire. This paper makes the following contributions.

\begin{itemize}
\item \textbf{A statement of the DNP3 timing-and-shape leak on real hardware.} We measure the native
cross-layer response time and segment shape of an SEL-751 relay and show that they separate its read
and select transaction classes.
\item \textbf{An in-network timing normalization.} We hold the acknowledgment and the response in
switch queues and release them at fixed offsets from the request, driving the master-visible CLRT to a
single policy value without a controller or an endpoint change.
\item \textbf{A byte-preserving shape normalization.} We replicate the eligible response and carve it
at an existing DNP3 CRC-block boundary into a fixed $[28,21]$ segment shape, so the master reassembles
a valid 49-byte response; we claim a CRC- and checksum-valid reconstruction, not byte identity to a
source frame.
\item \textbf{A control-command treatment and its scheduling.} We hold the OPERATE and pin the
master-visible acknowledgment and echo to the request arrival, and we show why the two treatments
require two independent switch scheduling lanes.
\item \textbf{A one-program implementation and a bounded hardware evaluation.} We implement the design
as one P4 program on a single switch pipeline and evaluate it on a physical SEL-751 with a
reproducible capture-to-claim analysis, an explicit safety result, and a stated evidence boundary.
\end{itemize}
